\begin{textsample}{POS Dim 2 – ro – Score 38.00 – ro/ro\_em\_92.txt}  \label{ex:f2_pos_181}
3.3.7.2.
Orientações para composição do Itinerário Formativo de EPT O \textbf{Currículo} do Ensino Médio, em sua reestruturação, possui a parte comum, denominada de Formação Geral, que considera as competências e habilidades definidas pela Base Nacional Comum Curricular – BNCC, e outra parte \textbf{flexível} que são os \textbf{Itinerários} Formativos, entendidos como um conjunto de unidades em que os estudantes podem escolher a partir do seu interesse para aprofundar e ampliar aprendizagens em uma ou mais Áreas de Conhecimento e/ou na Formação \textbf{Técnica} e Profissional.
Em relação a parte \textbf{flexível} do \textbf{currículo}, a \textbf{Portaria} no 1.432/2018/MEC, estabelece os Referenciais para a Elaboração dos \textbf{Itinerários} Formativos, sendo que, em seu anexo, consta a orientação aos sistemas de ensino para a construção dos \textbf{itinerários} formativos, de modo a \textbf{atender} o que preconizam a Lei no 13.415/2017 e as \textbf{Diretrizes} Curriculares Nacionais do Ensino Médio, \textbf{publicadas} na Resolução CNE/CEB N. 03/2018.
Nesta perspectiva, o Itinerário Formativo de Educação Profissional e Tecnológica ( EPT ), figura como uma oportunidade aos estudantes para que tenham um diferencial em sua formação.
Por este âmbito o mundo do trabalho passa a ser referência para a organização e \textbf{oferta} da Educação Profissional, visto que, dentre os desafios postos pela evolução tecnológica da sociedade moderna, está o conhecimento e domínio de processos de trabalho.
Nessa mesma linha de pensamento, no \textbf{Parecer} CNE/CEB no 11/2012 ( p. 11 ) é destacado que : > [... ] Educação Básica, formação profissional e aprendizagem ao longo da vida contribui significativamente para a promoção dos interesses individuais e coletivos dos \textbf{trabalhadores} e dos empregadores, bem como dos interesses sociais do desenvolvimento socioeconômico, especialmente, tendo em conta a importância fundamental do pleno emprego, da erradicação da pobreza, da inclusão social e do crescimento econômico sustentado.
Na concepção de integração do Ensino Médio com a EPT, como \textbf{política} pública, é necessário o envolvimento de todos os elementos constitutivos do Sistema Estadual de Ensino de Rondônia, bem como uma sincronizada articulação com os movimentos sociais, economias locais e sociedade civil em geral, por gerarem conhecimento em seus campos de atuação.
A formação integrada representa a superação do ser humano dividido historicamente pela dicotomia social do trabalho entre a ação de executar e a ação de pensar, dirigir ou planejar, projetando -o para além da redução da preparação para o trabalho ao seu aspecto operacional, simplificado, dissociado dos conhecimentos que estão na sua gênese científico-tecnológica e na sua apropriação histórico-social.
Como formação humana, será garantido ao \textbf{adolescente}, ao jovem e ao adulto \textbf{trabalhador} o direito a uma formação completa para a leitura do mundo e para a atuação como cidadão pertencente a um país, integrado dignamente à sua sociedade política, oferecendo -lhe “ [... ] a compreensão das relações sociais subjacentes a todos os fenômenos ” ( CIAVATTA, 2005, p. 85 ).
É importante destacar, de acordo com o documento denominado Educação Profissional \textbf{Técnica} de Nível Médio Integrada ao Ensino Médio ( 2007 ), que : > [... ] a história da humanidade é a história da produção da existência humana e a história do conhecimento é a história do processo de apropriação social dos potenciais da natureza para o próprio homem, mediada pelo trabalho.
Por isso, o trabalho é mediação ontológica e histórica na produção de conhecimento. ( BRASIL, 2007, p. 42 ).
Neste contexto, é necessário dizer que o ser humano é produtor de sua realidade e, por isso, apropria -se dela e pode transformá -la, lançando mão, para isso, do trabalho como prática econômica, para garantir a sua existência, produzindo riquezas e satisfazendo necessidades.
No Estado de Rondônia, a relação econômica vai se tornando fundamento da profissionalização, em decorrência do crescimento econômico a partir do agronegócio, da ampliação do mapa turístico, da pecuária, da piscicultura e da agricultura familiar, bem como as questões ambientais, que impactam no setor comercial e industrial da economia do Estado.
O Ensino Médio articulado com o \textbf{itinerário} de EPT representa a consonância de interesses dos estudantes, que buscam sua inserção no mundo do trabalho e no desenvolvimento do potencial socioeconômico, ambiental e regional.
Para tanto, é necessário articular o ambiente escolar com os familiares dos estudantes e a sociedade em geral, pois as experiências de formação integrada não se fazem no isolamento institucional.
A escola deve levar em conta a visão que os alunos têm de si mesmos, as possibilidades de inserção social e profissional que o mundo externo lhes oferece, tendo como ponto de partida as modalidades formativas oferecidas pela escola.
Isso exige um processo de diálogo e de conscientização dos estudantes e de suas famílias sobre as próprias expectativas e sua possível realização.
É necessário considerar ainda as necessidades materiais dos estudantes, bem como proporcionar condições didático-pedagógicas às escolas e aos professores, vez que o sistema de ensino e as escolas precisarão se adaptar de modo a \textbf{atender} as necessidades materiais para levar adiante este processo educacional.
O ponto primordial para tal empreitada situa -se na preparação da estrutura necessária, para comportar a proposta de articulação do Ensino Médio com o Itinerário de EPT, dotando -a dos recursos necessários conforme estrutura mínima estabelecida no \textbf{Catálogo} Nacional de \textbf{Cursos} \textbf{Técnicos} para cada \textbf{curso} \textbf{ofertado}, em consonância com os arranjos produtivos locais e outros recursos usados para o apontamento das demandas de \textbf{cursos} \textbf{técnicos}.
As escolas do Sistema de Ensino que optarem pelo \textbf{itinerário} de EPT de forma articulada com o Ensino Médio obterão a \textbf{regulamentação} \textbf{prevista} nas normas específicas em vigência para a \textbf{oferta} de \textbf{cursos} \textbf{técnicos}, de modo a permitir o protagonismo dos estudantes no constructo de seu perfil como profissional a se apresentar para o mundo do trabalho e de seu perfil como cidadão consciente e exitoso nas relações socioambientais em que estiver inserido.
As unidades de aprendizagem que integram a articulação do Ensino Médio com o \textbf{itinerário} de EPT consideram as demandas e necessidades do mundo contemporâneo, cuja prática pedagógica se desenvolve a partir da realidade social dos estudantes e dos aspectos teóricos que orientam as experiências de aprendizagem, constituindo a dimensão de uma intenção planejada e com visão ética sobre as percepções da vida.
Desta forma, as Unidades de Aprendizagem dialogam com os eixos estruturantes e com cada elemento da sequência didática, para que as escolas ofertantes do \textbf{itinerário} de EPT possam trabalhar os eixos com metodologias criadas de forma coletiva, proporcionando a aquisição de competências relacionadas ao mundo do trabalho.
Assim, a execução dos projetos relacionados aos eixos permite que os estudantes vivenciem atividades práticas de campo, de tal forma que apliquem os conhecimentos da parte propedêutica do \textbf{currículo} e compreendam o contexto socioeconômico e cultural em que estão inseridos, as relações do mundo do trabalho, e interpessoais existentes, os serviços comunitários etc.
Em consequência, a arquitetura curricular do Ensino Médio, considerando as novas formas de organização e de \textbf{gestão} do trabalho, que passam por modificações estruturais cada vez mais aprofundadas, requer um desenho que contemple uma Formação Geral Básica e Formação Profissional que dialoguem e estabeleçam as conexões necessárias.
Para tanto, é fundamental disponibilizar aos estudantes do Ensino Médio arranjos curriculares que lhe proporcionem a vivência de experiências que impactem em sua vida, de sua família e de sua comunidade.
Uma vez que o Itinerário Formativo de EPT assume, numa intrínseca relação com os demais \textbf{Itinerários} das Áreas de Conhecimento da Formação Geral Básica um papel importante na perspectiva de \textbf{atender} as necessidades e as expectativas dos jovens, contribuindo para que tenham maior interesse em estar na escola, sua permanência e, consequentemente, sua aprendizagem.
A implantação da Educação Profissional e Tecnológica como \textbf{itinerário} formativo efetiva -se a partir de algumas opções de composição dos \textbf{itinerários}, contemplando sempre tempos e espaços para a construção e o \textbf{fortalecimento} do Projeto de Vida e o desenvolvimento de competências e habilidades relacionadas à preparação geral para o mundo do trabalho.
Os percursos poderão contemplar também a \textbf{oferta} de \textbf{cursos} \textbf{técnicos}, de acordo com o CNCT, de \textbf{cursos} FIC organizados de forma articulada e os \textbf{cursos} de \textbf{qualificação} profissional.
O percurso formativo da preparação básica para o trabalho se organiza a partir da abordagem e da integração dos diferentes eixos estruturantes, sendo que as habilidades a eles associadas somam -se a outras habilidades básicas requeridas indistintamente pelo mundo do trabalho e as habilidades específicas requeridas pelas distintas ocupações.
Na \textbf{oferta} do \textbf{itinerário} de EPT, as escolas poderão aproveitar os conhecimentos resultantes da \textbf{oferta} de \textbf{cursos} de Formação Inicial e Continuada ( FIC ) e de \textbf{qualificação}, para \textbf{prosseguimento} de estudos, desde que diretamente relacionados com o perfil profissional de conclusão da respectiva \textbf{qualificação} profissional ou \textbf{habilitação} profissional \textbf{técnica} ou tecnológica, nos termos da \textbf{legislação} de ensino \textbf{vigente} e de acordo com os correspondentes planos de \textbf{curso}.
No contexto da articulação do Ensino Médio com o \textbf{itinerário} EPT, os \textbf{cursos} \textbf{técnicos} serão \textbf{ofertados} nas formas integrada e concomitante ao Ensino Médio, sendo integrada, quando \textbf{ofertada} somente a quem já tenha concluído o Ensino Fundamental, com matrícula única na mesma \textbf{instituição}, de modo a conduzir o estudante à \textbf{habilitação} profissional \textbf{técnica} ao mesmo tempo em que conclui a última etapa da Educação Básica ; a forma concomitante é \textbf{ofertada} a quem ingressa no Ensino Médio ou já estejam \textbf{cursando}, sendo necessárias matrículas distintas para cada \textbf{curso}, com oportunidades educacionais disponíveis, seja em unidades de ensino da mesma rede ou em distintas \textbf{instituições} e redes de ensino ; a forma concomitante pode ser intercomplementar, quando desenvolvida simultaneamente em distintas \textbf{instituições} ou redes de ensino, mas integrada no conteúdo, mediante a celebração de convênio ou acordo de intercomplementaridade para a execução de projeto pedagógico unificado.
A prática profissional supervisionada, \textbf{prevista} na organização curricular do \textbf{itinerário} de EPT, está relacionada aos correspondentes fundamentos \textbf{técnicos}, científicos e tecnológicos, orientada pelo trabalho como princípio educativo e pela pesquisa como princípio pedagógico, possibilitando ao estudante a preparação para o enfrentamento do desafio do desenvolvimento da aprendizagem permanente, integrando as \textbf{cargas} \textbf{horárias} mínimas de cada \textbf{habilitação} profissional \textbf{técnica} e tecnológica, uma vez que compreende diferentes situações de vivência profissional, aprendizagem e trabalho, a partir de experimentos e atividades específicas, desenvolvidas com diferentes recursos tecnológicos em oficinas, laboratórios ou salas ambientes na própria \textbf{instituição} de ensino ou em \textbf{entidade} parceira.
O estágio profissional supervisionado, quando \textbf{previsto} pela \textbf{instituição} em função do perfil de formação ou exigido pela natureza da ocupação, será incluído no correspondente plano de \textbf{curso}, em \textbf{conformidade} com a Lei n. 11.788, de 25 de setembro de 2008, que “ Dispõe sobre o estágio de estudantes [... ] e dá outras \textbf{providências} ”.
O estágio profissional supervisionado é desenvolvido em ambiente real de trabalho, como ato educativo e supervisionado pela \textbf{instituição} de ensino, em regime de parceria com organizações do mundo do trabalho, para a efetiva preparação do estudante para o trabalho, com plano de realização do estágio profissional supervisionado explicitando as correspondentes ações didáticas e pedagógicas, uma vez que é ato educativo de responsabilidade da escola ofertante do \textbf{itinerário} de EPT.
Os saberes adquiridos na EPT e no trabalho serão reconhecidos por meio de processo formal de avaliação e reconhecimento de saberes e competências profissionais, para fins de exercício profissional e de \textbf{prosseguimento} ou conclusão de estudos, conforme \textbf{disposto} no \textbf{artigo} 41, da Lei no 9.394/1996, abrangendo a avaliação do \textbf{itinerário} profissional e social do estudante, que inclui estudos não formais e experiência no trabalho e a orientação para continuidade de estudos.
A \textbf{certificação} será realizada por meio da emissão de certificados e diplomas de \textbf{cursos} de Educação Profissional e Tecnológica, para fins de exercício profissional e de \textbf{prosseguimento} e conclusão de estudos, cabendo às escolas adotar as \textbf{providências} para expedição e registro dos certificados e diplomas de \textbf{cursos} de Educação Profissional e Tecnológica sob sua responsabilidade no Sistema Nacional de Informações da Educação Profissional e Tecnológica ( SISTEC ).
Em suma, espera -se contribuir para a escola vencer o grande desafio da falta dos conhecimentos básicos e das novas condições requeridas para ingresso dos jovens no mundo do trabalho que são globalização dos meios de produção, do comércio e da indústria, e a utilização crescente de novas tecnologias, de modo especial, aquelas relacionadas com a informatização.
A superação dessa falha na formação dos estudantes do Ensino Médio bem como da Educação Profissional é fundamental para garantir seu desenvolvimento e sua cidadania.

% matched lemmas: adolescente, artigo, atender, carga, catálogo, certificação, conformidade, currículo, cursar, curso, diretriz, disposto, entidade, flexível, fortalecimento, gestão, habilitação, horário, instituição, itinerário, legislação, oferta, ofertar, parecer, política, portaria, prever, previsto, prosseguimento, providência, publicar, qualificação, regulamentação, trabalhador, técnico, vigente
\end{textsample}
